\chapter{Bảng hỏi khảo sát trải nghiệm người dùng}
\label{Appendix:survey}

Phụ lục này trình bày đầy đủ bộ câu hỏi của bảng khảo sát trải nghiệm người dùng đã sử dụng để thu thập đánh giá chủ quan ở Mục~\ref{sec:end_user_evaluation}. Bảng hỏi được triển khai trực tuyến bằng Google Biểu mẫu và người tham gia trả lời sau khi đã trực tiếp trải nghiệm hệ thống chạy thật tại \texttt{career-nova.online}. Kết quả tổng hợp tương ứng được phân tích trong Mục~\ref{sec:end_user_evaluation}; phần dưới đây chỉ liệt kê nội dung các câu hỏi theo đúng thứ tự bốn phần của bảng hỏi.

\section*{Phần I --- Thông tin phân loại người dùng}

\begin{enumerate}
    \item Bối cảnh học tập/nghề nghiệp hiện tại của bạn? \textit{(chọn một)}
    \begin{itemize}
        \item Sinh viên đang theo học
        \item Mới ra trường, đang tìm việc/thực tập
        \item Đã đi làm dưới một năm
        \item Đã đi làm trên một năm
    \end{itemize}
    \item Mục tiêu nghề nghiệp (vị trí) bạn đang hướng tới? \textit{(câu hỏi mở)}
    \item Tình trạng tìm việc hiện tại của bạn? \textit{(chọn một)}
    \begin{itemize}
        \item Đang tập trung trau dồi kỹ năng
        \item Đang tích cực ứng tuyển
        \item Đã đi làm
    \end{itemize}
\end{enumerate}

\section*{Phần II --- Giao diện \& Tốc độ}

\noindent Các phát biểu sau được đánh giá theo thang Likert 5 mức (1 --- Hoàn toàn không đồng ý; 5 --- Hoàn toàn đồng ý):

\begin{enumerate}
    \item Các thao tác trên ứng dụng thuận tiện, dễ dàng.
    \item Giao diện giúp tìm kiếm thông tin cần thiết nhanh chóng.
    \item Tốc độ phản hồi của trang web đáp ứng kỳ vọng.
    \item Bố cục thông tin trên Dashboard vừa vặn, không quá tải.
\end{enumerate}

\section*{Phần III --- Biểu đồ thị trường \& So khớp CV}

\noindent Các phát biểu đánh giá theo thang Likert 5 mức:

\begin{enumerate}
    \item Việc thiết lập và tải CV lên hệ thống diễn ra dễ dàng.
    \item Hệ thống nhận diện đầy đủ các kỹ năng liệt kê trong CV.
    \item Biểu đồ Radar giúp nắm bắt điểm mạnh/điểm yếu bản thân.
\end{enumerate}

\noindent Câu hỏi bổ sung:

\begin{enumerate}
    \setcounter{enumi}{3}
    \item Đánh giá mức độ hữu ích của từng biểu đồ thị trường sau (thang: Không hữu ích --- Ít hữu ích --- Hữu ích --- Rất hữu ích):
    \begin{itemize}
        \item Top 5 vị trí được tuyển nhiều nhất
        \item Top 10 kỹ năng được yêu cầu
        \item Xu hướng tuyển dụng theo thời gian
        \item Phân bố mức lương theo nhóm kỹ năng
    \end{itemize}
    \item Bạn mong muốn hệ thống bổ sung thêm loại thông tin thị trường nào? \textit{(chọn tối đa 2)}
    \begin{itemize}
        \item So sánh mặt bằng lương (Salary Benchmarking)
        \item Quy trình phỏng vấn \& bộ câu hỏi test kỹ thuật
        \item Xu hướng công nghệ mới nổi (Emerging Tech Trends)
        \item Phân tích Tech Stack/Dự án thực tế từ JD
    \end{itemize}
    \item Bạn cần thêm điều gì để tin tưởng con số ``Điểm phù hợp'' (Matching) hơn? \textit{(nhiều lựa chọn/câu hỏi mở, ví dụ: giải thích tiêu chí chấm điểm, hiển thị danh sách kỹ năng còn thiếu\ldots)}
\end{enumerate}

\section*{Phần IV --- Khóa học \& Lộ trình học tập}

\noindent Các phát biểu đánh giá theo thang Likert 5 mức:

\begin{enumerate}
    \item Các khóa học đề xuất sát với chuyên môn và nhu cầu thực tế.
    \item Lộ trình giúp biết cách bắt đầu cải thiện kỹ năng ngay.
    \item Lộ trình được cá nhân hóa theo mục tiêu nghề nghiệp.
    \item Tiết kiệm đáng kể thời gian tìm kiếm định hướng so với các công cụ khác.
\end{enumerate}

\section*{Câu hỏi tổng kết}

\begin{enumerate}
    \item Trên thang điểm 0--10, mức độ thiện cảm của bạn với sản phẩm (khả năng bạn sẽ tiếp tục sử dụng hoặc giới thiệu cho người khác)?
\end{enumerate}
